\section{Kebaruan Penelitian}
\label{sec:kebaruan-penelitian}

Kebaruan yang ditawarkan oleh penelitian ini terletak pada tiga aspek berikut:

\begin{enumerate}
    \item Spektrum frekuensi mode kuasinormal untuk perturbasi medan skalar bermassa pada ruang-waktu Schwarzschild Anti-de Sitter dihitung secara numerik, termasuk mode overtone pertama ($n=1$). Hasil ini memperbarui studi Horowitz dan Hubeny \parencite{Horowitz2000} yang hanya menghitung mode fundamental ($n=0$) untuk medan skalar tak bermassa pada latar belakang yang sama.

    \item Respons frekuensi mode kuasinormal terhadap variasi parameter yang sama, yaitu momentum sudut $\ell$, indeks overtone $n$, dan massa medan skalar $\mu$, dibandingkan secara langsung antara ruang-waktu Schwarzschild dan Schwarzschild Anti-de Sitter dalam satu studi. Perbandingan ini mengungkap bahwa variasi parameter yang sama dapat menghasilkan kecenderungan frekuensi yang berbeda arah pada kedua geometri, khususnya pada pengaruh massa medan skalar terhadap laju peluruhan perturbasi.

    \item Kestabilan kedua ruang-waktu terhadap perturbasi medan skalar, baik tak bermassa maupun bermassa, dianalisis secara komparatif. Meskipun kedua geometri sama-sama stabil, karakteristik peluruhannya berbeda: pada ruang-waktu Schwarzschild, peningkatan massa lubang hitam memperlambat peluruhan perturbasi, sedangkan pada ruang-waktu Schwarzschild Anti-de Sitter, peningkatan ukuran horizon peristiwa justru mempercepat peluruhan.
\end{enumerate}